% Chapter 6 (c6_similarity), chapter opening paragraph above the Introduction heading. Reason: repeats the Introduction's research questions (c6_similarity-P2). The sentence "A method can succeed at the first and fail at the third." was moved to the Methods paragraph on the three tasks.
The hybrid representation was first conceived to repair molecular similarity. This chapter tests whether features chosen for the task give better neighbors than single representations, and whether better neighbors help to find noisy labels, estimate their errors and train on them. A method can succeed at the first and fail at the third. Because the type of noise determines what can be corrected, we treat censoring, Gaussian noise and real labels separately.

% Chapter 6 (c6_similarity), Sample Removal Methods, opening sentence. Reason: restates the similar property principle stated in Chapter 2 (c6_similarity-P5); replaced by a back-reference to Chapter~\ref{chap:background}.
Distance-based approaches rely on the assumption that similar molecules have similar properties \citep{Cherkasov2014}.
